\chapter{Ethical Considerations and Data Management Plan}
\label{app:ethics}

\section{Ethical Considerations}

This study involves no human participants, no animal subjects and no personal
data. It is an instrumented agronomic experiment on a cultivated vegetable
crop, and accordingly it required no clearance from an ethics review board.
The ethical obligations that do apply to it are obligations of research
integrity, and they are stated here so that they can be checked.

% TODO: confirm with the advisers whether the department requires a formal
% ethics clearance form even for studies with no human or animal subjects.
% The predecessor study (Zamora et al., 2025) required one because it
% surveyed users; this study does not survey anyone.

\subsection{Integrity of the Record}

No measurement reported in this study was interpolated, back-filled or
reconstructed after the fact. Capture days lost to power interruption are
recorded as missing rather than filled from a neighbouring day, and the
affected days are named in Chapter~\ref{chap:results}. Every change to the
measurement geometry is logged with the date, the reason and a retained copy
of the previous configuration, so that any reported value can be traced to the
exact settings that produced it.

\subsection{Reporting of Unfavourable Results}

Results that weaken the study are reported alongside those that support it.
The failure of the treatment bands to separate as designed, the days lost from
the series, and the gap between the algorithm's agreement with human
annotation and the agreement humans reach among themselves are all stated in
Chapter~\ref{chap:results} rather than confined to a limitations section.

\section{Data Management Plan}

\subsection{Data Generated}

% TODO: confirm final figures at submission.
\begin{itemize}
  \item Daily overhead frames at 1296 by 1296 px, one per capture day.
  \item \texttt{growth\_dataset.csv}: one row per pot per capture day, holding
        the day index, timestamp, soil moisture, the several Green Pixel Ratio
        variants, pixel counts and quality flags.
  \item \texttt{sm\_log.csv}: periodic soil moisture snapshots for all nine
        pots.
  \item \texttt{irrigation\_log.csv}: one row per delivered dose.
  \item \texttt{power\_log.csv}, \texttt{error\_log.csv}: reliability records.
  \item \texttt{pot\_rois.json} and its dated predecessors: measurement
        geometry and its change history.
  \item \texttt{harvest\_weights.csv}: above-ground fresh weight per plant.
\end{itemize}

\subsection{Storage and Backup}

Working data are held on the controller's storage and mirrored to a version
controlled repository. Every change to the measurement configuration is
committed with the reason recorded in the file itself, so the history of the
instrument is recoverable alongside the history of the data.

\subsection{Retention and Access}

The dataset and the analysis code will be deposited with the Department of
Computer Engineering on completion of the study, so that the reported results
can be reproduced from the raw frames.
